%------------------------------------------------------------------------------%
\documentclass[fleqn,utf8,aspectratio=169,12pt]{beamer}

\usepackage{calculo_varias_variaveis-1}

\title{Curvas Paramétricas}

%------------------------------------------------------------------------------%
\begin{document}

\addtitle

%------------------------------------------------------------------------------%
\section{Introdução}

\framedgraphic{Trajetória}{trajeto_maps}{}

%------------------------------------------------------------------------------%
\begin{frame}{Curvas Paramétricas}
  Funções de $\R$ em $\R^n$
  \[
    (x, y) = \big( f(t), g(t) \big)
  \]
  ou
  \[
  \left[\begin{array}{>{\,}c<{\,}}
    x \\[1mm]
    y
  \end{array}\right]
  =
  \left[\begin{array}{>{\,}c<{\,}}
    f(t) \\[1mm]
    g(t)
  \end{array}\right]
  \]

  \pause
  \vspace{\baselineskip}
  Podemos pensar na variável como o tempo e a curva como uma trajetória
\end{frame}

%------------------------------------------------------------------------------%
\section{Curvas Paramétricas}

%------------------------------------------------------------------------------%
\begin{frame}{Definição em 2D}
  Se $x$ e $y$ são dados como funções
  \[
    x = f(t)
    \qquad
    y = g(t)
  \]
  sobre um intervalo $I$ de valores de $t$,
  \pause
  então o conjunto de pontos
  \[
    (x, y) = \big( f(t), g(t)\big)
  \]
  \pause
  forma uma \emph{curva paramétrica} em $\R^2$

  \pause
  As equações são as \emph{equações paramétricas} da curva
\end{frame}

%------------------------------------------------------------------------------%
\begin{frame}{Definição em 3D}
  Se $x$, $y$ e $z$ são dados como funções
  \[
    x = f(t)
    \qquad
    y = g(t)
    \qquad
    z = h(t)
  \]
  sobre um intervalo $I$ de valores de $t$,
  então o conjunto de pontos
  \[
    (x, y, z) = \big( f(t), g(t), h(t) \big)
  \]
  forma uma \emph{curva paramétrica} em $\R^3$
\end{frame}

%------------------------------------------------------------------------------%
\begin{frame}{Definição Geral}
  Se $x_i$, $i=1, \dots, n$ são as coordenadas no espaço $\R^n$
  \[
    x_i = f_i(t)
  \]
  sobre um intervalo $I$ de valores de $t$,
  então o conjunto de pontos
  \[
    (x_1, x_2, \dots, x_n) = \big( f_1(t), f_2(t), \dots, f_n(t)\big)
  \]
  forma uma \emph{curva paramétrica} em $\R^n$
\end{frame}

%------------------------------------------------------------------------------%
\begin{frame}{Nomenclatura}
  \emph{$t$} \quad é o parâmetro da curva

  \pause
  \emph{$I$} \quad intervalo do parâmetro

  \pause
  \vspace{\baselineskip}
  Se $I$ é um intervalo fechado \emph{\(a\leq t \leq b\)}
  \pause
  \begin{description}[<+->]
    \item[\(\big(f(a), g(a)\big)\)] é o ponto inicial
    \item[\(\big(f(b), g(b)\big)\)] é o ponto final
  \end{description}
\end{frame}

%------------------------------------------------------------------------------%
\begin{frame}{Notações}
  \[
    (x, y)
    = \big(x(t), y(t)\big)
    = \big(f(t), g(t)\big)
  \]

  \pause
  \[
      \begin{pmatrix}  x    \\ y    \end{pmatrix}
    = \begin{pmatrix}  x(t) \\ y(t) \end{pmatrix}
    = \begin{pmatrix}  f(t) \\ g(t) \end{pmatrix}
  \]

  \pause
  \[
    (x_1, x_2, \dots, x_n) = \big( f_1(t), f_2(t), \dots, f_n(t)\big)
  \]
\end{frame}

%------------------------------------------------------------------------------%
\section{Exemplos}

\include{B-02-Curvas_parametricas-ex}

%------------------------------------------------------------------------------%
\section{Lista Mínima}

%------------------------------------------------------------------------------%
\begin{frame}{Lista Mínima}
  Cálculo Vol. 2 do Thomas 12\textsuperscript{a} ed. -- Seção 11.1
  \begin{enumerate}
    \item Estudar todo o texto da seção
    \item Resolver os exercícios: 3, 5, 11, 19, 21, 23, 31
  \end{enumerate}

  \vfill
  Atenção: A prova é baseada no livro, não nas apresentações
\end{frame}

%------------------------------------------------------------------------------%
\end{document}
%------------------------------------------------------------------------------%
